\pagenumbering{roman}
\chapter*{Dédicace}
\addcontentsline{toc}{chapter}{Dédicace}
\markboth{Dédicace}{Dédicace}
\vspace*{3cm}
\begin{center}\Large\color{navy}\emph{(à compléter)}\end{center}

\chapter*{Remerciements}
\addcontentsline{toc}{chapter}{Remerciements}
\markboth{Remerciements}{Remerciements}
\begin{flushright}\itshape
« Une seule main ne peut nouer un fagot de bois. »\\
\upshape\small --- Adage africain
\end{flushright}
\vspace{0.5em}
Fort de cet adage, je prends la plume pour remercier celles et ceux dont le soutien m'a permis de mener à terme la rédaction de ce travail. Ce travail n'aurait pas été possible sans l'accompagnement de plusieurs personnes, à qui j'exprime ma sincère gratitude :
\begin{itemize}
\item \textbf{Dr Hervé Talé Kalachi}, mon encadrant, pour ses relectures exigeantes et les consignes méthodologiques qui ont structuré ce travail ;
\item \emph{(à compléter : membres du jury, enseignants du Département de Génie Informatique, collègues, famille, amis.)}
\end{itemize}

\chapter*{Abstract}
\addcontentsline{toc}{chapter}{Abstract}
MAYO is an Oil-and-Vinegar multivariate signature scheme submitted to post-quantum standardisation. In the official AVX2 C implementation, products of the public matrix $P_1$ by the vectors of $V$ are computed with \code{vpshufb} lookups on packed $\mathbb{F}_{16}$ vectors. When the vector length $m$ is not a multiple of 64, the official kernel covers the tail with an overlapping 32-byte load that duplicates work. We study a local change, called A2, which handles the 8-byte tail of each $P_1$ entry in a single 128-bit register and one xmm \code{vpshufb} on the low lane of the multiplication table, reducing the number of shuffles per scalar pair from 4 to 3 (MAYO-1) and from 6 to 5 (MAYO-5). We formalise the counting argument and its limits, establish the properties on which correctness rests, and measure the effect on a physical laptop with a paired protocol. Kernel gains are 12--29\,\%; the full-signature gain is 10.3\,\% for MAYO-1 but only 1.3\,\% for MAYO-5. Static analysis shows that in the signature context GCC~16.2.0 emits extra 32-byte stores in several loops, asymmetrically between the two versions, consistently with the measured loss; causality is not established. GCC's dumps show that the extra stores are per-iteration stores of the accumulator arrays that are not moved out of the loop; SRA is ruled out, and why the compiler makes this choice remains open.


\chapter*{Résumé}
\addcontentsline{toc}{chapter}{Résumé}
MAYO est un schéma de signature numérique « multivarié » de type Oil-and-Vinegar. Il est candidat à la standardisation de la cryptographie post-quantique. Dans le code C officiel (MAYO-C) pour AVX2, le programme multiplie la matrice publique $P_1$ par les vecteurs de $V$ avec l'instruction \code{vpshufb}. Les vecteurs de $\mathbb{F}_{16}$ y sont rangés à deux éléments par octet. Quand la longueur $m$ n'est pas un multiple de 64, le code officiel traite la fin du vecteur en relisant 32 octets déjà traités. Cela fait des calculs en double : 48 éléments sur 128 pour MAYO-1 ($m=80$) et 50 sur 192 pour MAYO-5 ($m=142$).

Notre travail étudie une petite modification, appelée A2. Elle traite la fin de 8 octets de chaque entrée de $P_1$ dans un seul registre de 128 bits : un seul \code{vpshufb} xmm remplace deux \code{vpshufb} ymm. Le nombre de \code{vpshufb} par paire de nombres passe de 4 à 3 (MAYO-1) et de 6 à 5 (MAYO-5). Nous démontrons cette limite de comptage, nous établissons les propriétés qui rendent A2 correct, puis nous mesurons l'effet sur un ordinateur portable (Intel Core i9-9880H, GCC~16.2.0). Les mesures comparent les deux versions dans les mêmes conditions, et l'unité statistique est le processus (une exécution du programme).

Voici les résultats établis. Sur les noyaux pris seuls, A2 réduit le temps de 19 à 22\,\% (\code{P1_times_Vt}) et de 12 à 29\,\% (\code{P1t_times_Vt}). Sur la signature complète, le gain est de 10,3\,\% pour MAYO-1, mais de seulement 1,3\,\% pour MAYO-5. Pour MAYO-5, le résultat principal est donc 1,3\,\% ; les 7,2\,\% obtenus avec le binaire instrumenté (toutes les sondes actives) servent seulement à expliquer pourquoi le protocole de mesure compte. En lisant le code produit par GCC~16.2.0 (sans l'exécuter), nous voyons que, sans toutes les sondes, le compilateur ajoute dans plusieurs boucles une écriture mémoire par accumulateur (32 octets pour un registre ymm, 16 octets pour un xmm), et que ce n'est pas la même chose dans les deux versions. Les temps mesurés varient de la même façon, mais le lien de cause à effet n'est pas démontré, et un modèle à un seul paramètre exagère le gain d'environ 40\,\%. Ces écritures en plus sont les stockages des tableaux d'accumulateurs, refaits à chaque tour de boucle au lieu d'être faits une seule fois après la boucle. Nous écartons la passe SRA du compilateur comme cause. Pourquoi le compilateur fait ce choix reste une question ouverte.

\medskip
\noindent\textbf{Mots-clés :} MAYO, signatures multivariées, cryptographie post-quantique, AVX2, \code{vpshufb}, $\mathbb{F}_{16}$, optimisation, mesure de performance.


\chapter*{Liste des abréviations}
\addcontentsline{toc}{chapter}{Liste des abréviations}
\markboth{Liste des abréviations}{Liste des abréviations}
\begin{itemize}[leftmargin=1.2cm,itemsep=2pt]
\item[\textbf{AVX2 :}] \emph{Advanced Vector Extensions 2} : instructions des processeurs qui traitent 256 bits (32 octets) à la fois
\item[\textbf{CT :}] \emph{Constant Time} (temps constant)
\item[\textbf{DSE :}] \emph{Dead Store Elimination} (passe de GCC)
\item[\textbf{FRE :}] \emph{Full Redundancy Elimination} (passe de GCC)
\item[\textbf{GCC :}] \emph{GNU Compiler Collection}
\item[\textbf{GFNI :}] \emph{Galois Field New Instructions}
\item[\textbf{IC95 :}] intervalle de confiance à 95\,\%
\item[\textbf{KAT :}] \emph{Known-Answer Test} : test de « réponses connues » (on compare à des valeurs officielles)
\item[\textbf{LIM :}] \emph{Loop Invariant Motion} (passe de GCC)
\item[\textbf{MAYO :}] schéma de signature multivariée de type Oil-and-Vinegar
\item[\textbf{FIPS :}] \emph{Federal Information Processing Standards} : normes fédérales américaines de traitement de l'information
\item[\textbf{MQ :}] \emph{Multivariate Quadratic} : équations de degré 2 à plusieurs variables
\item[\textbf{NIST :}] \emph{National Institute of Standards and Technology} : Institut national américain des normes et de la technologie
\item[\textbf{OV, UOV :}] \emph{(Unbalanced) Oil and Vinegar}
\item[\textbf{PRE :}] \emph{Partial Redundancy Elimination} (passe de GCC)
\item[\textbf{SIMD :}] \emph{Single Instruction, Multiple Data}
\item[\textbf{SRA :}] \emph{Scalar Replacement of Aggregates} (passe de GCC)
\item[\textbf{TSC :}] \emph{Time Stamp Counter} : compteur de ticks du processeur
\end{itemize}

\chapter*{Glossaire}
\addcontentsline{toc}{chapter}{Glossaire}
\markboth{Glossaire}{Glossaire}
\begin{itemize}[leftmargin=0.5cm,itemsep=3pt]
\item \textbf{$m$-vec :} vecteur de $m$ éléments de $\mathbb{F}_{16}$, empaqueté à deux éléments par octet.
\item \textbf{Limbe :} mot de 64 bits d'un accumulateur.
\item \textbf{Voie (\emph{lane}) :} moitié de 128 bits d'un registre ymm.
\item \textbf{Table de multiplication (multab) :} table de 16 octets donnant, pour une paire de nombres, les deux produits par un quartet.
\item \textbf{Corps fini :} ensemble fini de nombres où l'on peut additionner, soustraire, multiplier et diviser (par exemple $\mathbb{F}_{16}$).\item \textbf{Forme quadratique :} expression de degré 2 en plusieurs variables, par exemple $x_1x_2+x_3^2$.\item \textbf{Empreinte (hachage) :} courte suite de nombres calculée à partir d'un message ; deux messages différents ont des empreintes différentes.\item \textbf{Sel :} valeur aléatoire ajoutée au message avant le calcul de l'empreinte.\item \textbf{Compilateur, assembleur :} le compilateur transforme le code C en assembleur, la suite d'instructions que le processeur exécute.\item \textbf{Passe (du compilateur) :} une étape d'optimisation du code (par exemple SRA, PRE, DSE).\item \textbf{Bootstrap :} méthode qui rejoue le calcul sur des tirages au hasard des mesures pour estimer son incertitude.\item \textbf{Fin de vecteur :} fin d'un $m$-vec qui ne remplit pas un registre complet de 32 octets.
\item \textbf{Recouvrement :} chargement reculé du dernier registre, qui retraite des octets déjà traités.
\item \textbf{SAFE :} variante d'A2 qui ne lit rien au-delà de l'entrée de $P_1$.
\item \textbf{Tick TSC :} unité du compteur d'horodatage (\code{rdtscp}) ; ce n'est pas un cycle cœur.
\item \textbf{Processus :} unité statistique : une exécution du banc, 3\,001 paires.
\item \textbf{Doublon :} écritures de 32 octets supplémentaires par accumulateur dans le code généré, selon la construction.
\item \textbf{E1, E1b :} analyses statiques du code généré (E1b : origine du doublon, exécutée le 25 octobre 2026 sous GCC~16.2.0).
\item \textbf{Fait, inférence, hypothèse :} statuts des énoncés, notés \F{}, \I{}, \Hy{}.
\end{itemize}
